\chapter{Introduzione}
\label{cap:introduzione}

\todoscrivi{Scaletta consigliata. Ogni sezione risponde a una domanda precisa del lettore.}

\section{Contesto e motivazione}
\label{sec:contesto}
% Perche' il tema e' rilevante ORA. Dal generale (IA generativa nell'arte) allo
% specifico (il problema della creativita' macchinica). Evitare l'incipit
% enciclopedico: entrare nel merito entro la prima pagina.

\section{Definizione del problema}
\label{sec:problema}
% Qual e' il problema aperto, e perche' le soluzioni esistenti non lo risolvono.
% Questa sezione deve rendere inevitabile la domanda di ricerca che segue.

\section{Domande di ricerca}
\label{sec:domande-ricerca}
% ADR-0007 (2026-09-14): la RQ non e' piu' empirica (non chiede piu' se un
% meccanismo aumenta l'ambiguita' misurabile), e' storico-critica: come si e'
% evoluta, da Elgammal et al. (2017) a oggi, l'argomentazione contro la
% rivendicazione di creativita' dei sistemi generativi, e quali obiezioni
% restano aperte. Una sola domanda, verificabile sulla rassegna del cap. 3.

\section{Obiettivi e contributi}
\label{sec:obiettivi}
% ADR-0007: il contributo e' una rassegna critica organizzata cronologicamente,
% non un nuovo esperimento. Dichiarare esplicitamente il limite: nessun
% esperimento originale in questo documento (il lavoro DCGAN/CAN/E8 e' escluso
% per scelta, non per mancanza -- vedi D-028 in registro-decisioni.md).

\section{Struttura della tesi}
\label{sec:struttura}
% Un paragrafo per capitolo. Da scrivere alla fine.
